\newgeometry{hmargin=16mm,vmargin=19mm}
\section{Wolfram-Eingaben und Originalausgaben zur Fassung 0.1.0}
\label{app:wolfram-output}

Die Fenster zeigen links das Eingabeskript und rechts die zugehörige Originalausgabe. Die Zeilenangaben beziehen sich auf die jeweiligen Dateien. Die Blöcke eines Laufs setzen den vorangehenden Wolfram-Zustand voraus. Lange Zeilen werden für den Satz umgebrochen.

\medskip
\noindent
\begin{tabular}{@{}lp{.59\linewidth}r@{}}
\hline
\textbf{Lauf} & \textbf{Rechnungen} & \textbf{Seite} \\
\hline
HC und Feldalgebra & Projektionen, Biquaternionen und geometrische Identitäten & \pageref{app:wl-core} \\
Dixon-Träger & Kommutanten und Hyperladung & \pageref{app:wl-dixon} \\
Koide & Wurzeltripel, Gram-Spektrum und Massenauslese & \pageref{app:wl-koide} \\
OP5 & Quellen, Achsenprofil und Vergleichsrechnungen & \pageref{app:wl-op5} \\
\hline
\end{tabular}

\newcommand{\wolframdatei}[3]{%
  \ifnum\pdfstrcmp{#1}{dixon/dixon.wl}=0
    \VerbatimInput[firstline=#2,lastline=#3,fontsize=\footnotesize,baselinestretch=.88,%
      breaklines,breakafter={,[-/}]{audit/wolfram-v0.1.0/#1}%
  \else
    \VerbatimInput[firstline=#2,lastline=#3,fontsize=\footnotesize,baselinestretch=.88,%
      breaklines,breakanywhere]{audit/wolfram-v0.1.0/#1}%
  \fi}
\newcommand{\wolframfenster}[7]{%
  \subsubsection*{#1}%
  \noindent
  \begin{tabular}{|p{\dimexpr.5\linewidth-2\tabcolsep-1.5\arrayrulewidth\relax}%
                  |p{\dimexpr.5\linewidth-2\tabcolsep-1.5\arrayrulewidth\relax}|}
  \hline
  \rule{0pt}{2.5ex}\textbf{Wolfram-Input} \hfill \footnotesize Zeilen #3--#4 &
  \rule{0pt}{2.5ex}\textbf{Originalausgabe} \hfill \footnotesize Zeilen #6--#7 \\[.35ex]
  \hline
  \wolframdatei{#2}{#3}{#4}\vspace{.7\baselineskip} &
  \wolframdatei{#5}{#6}{#7}\vspace{.7\baselineskip} \\
  \hline
  \end{tabular}\par}

\medskip
\subsection{Heisenberg-Kompensator und Feldalgebra}\label{app:wl-core}
\wolframfenster{Projektion und ebene Welle}{core/article_core.wl}{1}{30}{core/article_core.txt}{1}{4}

\clearpage
\wolframfenster{Feldinvarianten, Biquaternionen und Nullspinoren}{core/article_core.wl}{31}{59}{core/article_core.txt}{5}{11}

\clearpage
\wolframfenster{Abschlussbedingung und Rang-eins-Vergleich}{core/article_core.wl}{60}{76}{core/article_core.txt}{12}{14}

\clearpage
\wolframfenster{Auslese, Eichung, Quellenreduktion und Koframe}{core/article_core.wl}{77}{102}{core/article_core.txt}{15}{23}

\clearpage
\subsection{Dixon-Träger und Kommutanten}\label{app:wl-dixon}
\wolframfenster{Multiplikation, Oktonionenderivationen und Farbkommutant}{dixon/dixon.wl}{1}{35}{dixon/dixon.txt}{1}{3}

\clearpage
\wolframfenster{Schwache Generatoren und Hyperladung}{dixon/dixon.wl}{36}{50}{dixon/dixon.txt}{4}{6}

\clearpage
\wolframfenster{Linearer Kommutant, Spektrum und antilineare Bedingung}{dixon/dixon.wl}{51}{77}{dixon/dixon.txt}{7}{16}

\clearpage
\subsection{Koide und Massenauslese}\label{app:wl-koide}
\wolframfenster{Zerlegung des Wurzeltripels und Phasenbereich}{koide/koide.wl}{1}{21}{koide/koide.txt}{1}{7}

\clearpage
\wolframfenster{HC-Projektion und Koisometrie}{koide/koide.wl}{22}{31}{koide/koide.txt}{8}{11}

\clearpage
\wolframfenster{Gram-Spektrum, Bargmann-Phase und Oktant}{koide/koide.wl}{32}{51}{koide/koide.txt}{12}{21}

\clearpage
\wolframfenster{Mischung, Interferenz und Massenskala}{koide/koide.wl}{52}{68}{koide/koide.txt}{22}{34}

\clearpage
\subsection{Geometrischer Vergleich in OP5}\label{app:wl-op5}
\wolframfenster{Quellterme und gekoppelte elastische Gleichungen}{op5/op5.wl}{1}{29}{op5/op5-output.txt}{1}{8}

\clearpage
\wolframfenster{Lineare Bilanz, Quellkontraktionen und Feldantwort}{op5/op5.wl}{30}{50}{op5/op5-output.txt}{9}{18}

\clearpage
\wolframfenster{Transversales Profil, Achsenkontrast und Kernintegral}{op5/op5.wl}{51}{65}{op5/op5-output.txt}{19}{24}

\clearpage
\wolframfenster{Quadratische Terme und Kontrollen des Quotienten}{op5/op5.wl}{66}{87}{op5/op5-output.txt}{25}{36}
